\documentclass[answers]{exam}
\usepackage{../../mypackages}
\usepackage{../../macros}

\usepackage{ltablex}
\keepXColumns


\title{Barème de notation des cahiers}
\author{N. Bancel}
\date{2 Décembre 2025}

% Mise en forme des solutions en bleu
\SolutionEmphasis{\color{blue}}
\renewcommand{\solutiontitle}{\noindent}

\definecolor{lavande}{RGB}{180,190,255}

\begin{document}

\textbf{Collège Lycée Suger}
\hfill
\textbf{Mathématiques} \\

\textbf{Année 2025-2026}
\hfill
\textbf{Classe de 6ème} \par

{\let\newpage\relax\maketitle}

%\begin{center}
%  \textbf{\textcolor{blue}{Liste des attendus pour la tenue du cahier d'exercices et du lutin}}
%\end{center}

\vspace{0.5cm}

\renewcommand{\arraystretch}{1.7}

\newcolumntype{Y}{>{\hsize=0.7\hsize}X} % Colonne plus étroite
\newcolumntype{Z}{>{\hsize=1.3\hsize}X} % Colonne plus large
\newcolumntype{L}[1]{>{\raggedright\arraybackslash}p{#1}}

% Nouveau tableau : colonnes larges pour Critère + Description
\centering
\begin{adjustbox}{width=1.2\textwidth, center}
\begin{longtable}{|L{0.17\textwidth}|L{0.23\textwidth}|L{0.48\textwidth}|c|}
\rowcolor{lavande}
\hline
\textbf{Catégorie} & \textbf{Critère évalué} & \textbf{Description} & \textbf{Points} \\[2pt]

% ---------------- Cahier d'exercices ----------------
\hline
\multirow{3}{*}{Cahier d’exercices}
  & Titres des exercices présents et soignés (+ utilisation de couleur)
  & Les exercices comportent un titre clair, souligné ou encadré. Et une référence claire à l'endroit où retrouver l'exercice ("Exercice N°34 page 127 - Le livre scolaire" ou "Exercice d'application N°3 - Lutin" ou "Exercice N°2 - Fiche d'exercice")
  & 2 \\
  \cline{2-4}

  & Ordres de grandeurs notés
  & Les ordres de grandeur sont indiqués lorsque nécessaire (unités, puissances de 10, cohérence des résultats).
  & 1 \\
  \cline{2-4}

  & Corrections notées (ou rattrapées en cas d'absence)
  & Toutes les corrections sont recopiées, complétées ou rattrapées en cas d’absence, de manière lisible.
  & 3 \\[6pt]
  

% ---------------- Lutin ----------------
\hline
\multirow{5}{*}{Lutin}
  & Feuilles rangées dans l’ordre
  & Les feuilles sont classées dans l’ordre attendu : leçons, exercices, évaluations…
  & 1 \\
  \cline{2-4}

  & Chaque feuille imprimée dans un feuillet transparent (1 page par feuillet)
  & Chaque feuille est protégée dans une pochette transparente, une seule par poche.
  & 1 \\
  \cline{2-4}

  & Interrogations + calcul mental faciles à retrouver
  & Les évaluations et exercices de calcul mental sont regroupés et rapidement identifiables.
  & 2 \\
  \cline{2-4}

  & Cours complet (rien de manquant)
  & Aucune feuille manquante : l’ensemble du cours est présent dans le lutin.
  & 2 \\
  \cline{2-4}

  & Leçons, phrases, définitions recopiées correctement
  & Le travail est propre : tracés à la règle ou compas, titres alignés, soin général.
  & 2 \\[6pt]

% ---------------- Général ----------------
\hline
\multirow{5}{*}{Général}
  & Organisation générale (titres soulignés, couleurs…)
  & Le cahier et le lutin sont organisés : titres visibles, couleur raisonnée, structure claire.
  & 2 \\
  \cline{2-4}

  & Écriture soignée et lisible
  & L’écriture est propre, lisible, régulière.
  & 1 \\
  \cline{2-4}

  & Peu ou pas de ratures
  & Les ratures sont limitées et propres, sans nuire à la lecture.
  & 1 \\
  \cline{2-4}

  & Pas de feuilles volantes
  & Aucune feuille glissée en vrac ; tout est rangé ou collé.
  & 1 \\
  \cline{2-4}

  & Présence de prénom / nom // Identification facile du lutin
  & Nom, prénom et classe sont visibles clairement.
  & 1 \\[6pt]

% ---------------- Total ----------------
\hline
\rowcolor{green!20}
 & \textbf{Total} & \textbf{Note maximale possible} & \textbf{20} \\
 \hline

\end{longtable}
\end{adjustbox}

\end{document}
